% !LW recipe = make-deck
% Level 8 (after the tracks meet): an optimiser. Theory: query trees
% rewritten by laws; what a correct optimiser is. Lean: a goal theorem.
% Build: make (rules in ../common/deck.mk).
\DocumentMetadata{
  lang        = en-GB,
  pdfstandard = ua-2,
  tagging     = on,
}
\documentclass{ltx-talk}

\input{preamble}

\title{Level 8: An Optimiser}
\author{The Proofs and Programs Club}
\institute{Royal Holloway, University of London}

\begin{document}

\leveltitle{8}{An Optimiser}{After the tracks meet: leans on TPIL}
  {After \levellink{7}}{Next: \levellink{9}, if you have not done it}

\levelframe{8}{Level 8 at a Glance}{%
  \code{optimise}, a function from queries to queries.}{%
  that \code{optimise q} and \code{q} have the same answer, on every
  database.}{%
  write recursive functions from queries to queries &
    \rds{tpil-structrec} \\
  prove by induction on a query, one case for each constructor &
    \rds{tpil-indtac, fpil-indother} \\
  write proofs that follow your function's cases & \rds{tpil-funind} \\
  prove termination when the recursion is not structural &
    \rds{tpil-wf} \\
}

% ---------------------------------------------------------------- theory
\theorypart{Query trees, rewritten by the laws}
\begin{frame}
  \frametitle{An Optimiser Rewrites Query Trees by the Laws}
  {\centering
    \begin{tikzpicture}[alt = {Two query trees with the same answer, where
        p is dept equals CS and q is topic equals Lists. Before: project
        name, over select p and q, over the join of (Member join Attended)
        with Session. After laws 1, 6 and 10: project name, over the join
        of (select p of Member, joined with Attended) with (select q of
        Session). The selections now sit next to the tables.},
      every node/.style = {inner sep = 2pt},
      tbl/.style = {font = \itshape},
      level distance = 9mm,
      level 2/.style = {sibling distance = 32mm},
      level 3/.style = {sibling distance = 22mm},
      level 4/.style = {sibling distance = 20mm},
      edge from parent/.style = {draw, black!50, thick}]
      \node at (0, 0) {$\proj{\att{name}}$}
        child {node {$\sel{p \wedge q}$}
          child {node {$\join$}
            child {node {$\join$}
              child {node[tbl] {Member}}
              child {node[tbl] {Attended}}}
            child {node[tbl] {Session}}}};
      \node at (9.6, 0) {$\proj{\att{name}}$}
        child {node {$\join$}
          child {node {$\join$}
            child {node {$\sel{p}$}
              child {node[tbl] {Member}}}
            child {node[tbl] {Attended}}}
          child {node {$\sel{q}$}
            child {node[tbl] {Session}}}};
      \draw[->, very thick, accent] (2.4, -1.8) -- node[above, note]
        {laws 1, 6, 10} (5.8, -1.8);
    \end{tikzpicture}\par}
  \medskip
  With $p$: $\att{dept} = \val{CS}$ and $q$: $\att{topic} = \val{Lists}$.
  A query is a term (Definition 5F.3), drawn as a tree. Same answer, but
  only CS members and the Lists session reach the joins.
  \theorytag{8}
\end{frame}

\begin{frame}
  \frametitle{A Correct Optimiser Never Changes an Answer}
  \defn{Definition 8.1 (Correct optimiser)}{An \emph{optimiser} is a
    function $\mathit{opt}$ from queries to queries. It is \emph{correct}
    if $\lBrack \mathit{opt}(q) \rBrack = \lBrack q \rBrack$
    for every query $q$ and every database (Definition 5F.4).}
  \begin{itemize}
    \item To prove it: induction on $q$, one case for each operator, as on
      lists in Levels 5T and 6.
    \item In each case, the rewrite is one of the laws of Level 7, and the
      law's side condition is exactly what $\mathit{opt}$ must check before
      it rewrites.
  \end{itemize}
  \theorytag{8}
\end{frame}

% ---------------------------------------------------------------- Lean
\leanpart{A correct optimiser}
\begin{frame*}
  \frametitle{Definition 8.1 in Lean Is One Theorem}
  \leanblock{goal.lean}
  \sketch
  Prove it by induction on \code{q}, using your laws of Level 7. Start
  small: an optimiser that only merges two selections into one (law 1).
  \leantag{8}
\end{frame*}

\end{document}
